\documentclass[11pt,a4paper]{article}

\usepackage[utf8]{inputenc}
\usepackage[T1]{fontenc}
\usepackage{lmodern}
\usepackage[french]{babel}

\usepackage[margin=2.5cm]{geometry}
\usepackage{amsmath,amssymb,amsthm}
\usepackage{graphicx}
\graphicspath{{figures/}}
\usepackage{xcolor}
\usepackage{booktabs}
\usepackage{hyperref}
\usepackage{tcolorbox}
\tcbuselibrary{theorems,breakable,skins}
\usepackage{enumitem}
\usepackage{fancyhdr}
\usepackage{array}
\usepackage{pifont}
\usepackage{listings}
\usepackage{multirow}

\newcommand{\cmark}{\ding{51}}
\newcommand{\xmark}{\ding{55}}
\newcommand{\bigO}{\mathcal{O}}

\hypersetup{
    colorlinks=true,
    linkcolor=blue!60!black,
    urlcolor=blue!60!black,
    citecolor=blue!60!black
}

\newtcolorbox{objectifbox}[1][]{colback=blue!5,colframe=blue!50,
    fonttitle=\bfseries,title=#1,breakable}
\newtcolorbox{infobox}[1][]{colback=green!5,colframe=green!50!black,
    fonttitle=\bfseries,title=#1,breakable}
\newtcolorbox{attentionbox}[1][]{colback=orange!5,colframe=orange!70,
    fonttitle=\bfseries,title=#1,breakable}
\newtcolorbox{retenirbox}[1][]{colback=yellow!10,colframe=yellow!50!black,
    fonttitle=\bfseries,title=#1,breakable}
\newtcolorbox{prereqbox}[1][]{colback=cyan!5,colframe=cyan!50!black,
    fonttitle=\bfseries,title=#1,breakable}
\newtcolorbox{theorembox}[1][]{colback=purple!5,colframe=purple!50,
    fonttitle=\bfseries,title=#1,breakable}
\newtcolorbox{appliboxQ}[1][]{colback=gray!8,colframe=gray!60,
    fonttitle=\bfseries,title=#1,breakable}
\newtcolorbox{appliboxR}[1][]{colback=green!3,colframe=green!40!black,
    fonttitle=\bfseries\itshape,title=#1,breakable}

\lstdefinestyle{java}{
  language=Java,
  basicstyle=\ttfamily\small,
  backgroundcolor=\color{gray!8},
  frame=single, framerule=0.5pt, rulecolor=\color{gray!40},
  keywordstyle=\color{blue!70!black}\bfseries,
  commentstyle=\color{green!50!black}\itshape,
  stringstyle=\color{red!60!black},
  numbers=left, numberstyle=\tiny\color{gray},
  breaklines=true, showstringspaces=false, tabsize=4,
  xleftmargin=4pt, xrightmargin=4pt, aboveskip=8pt, belowskip=8pt,
  escapeinside={(@}{@)},
  literate={é}{{\'e}}1 {è}{{\`e}}1 {ê}{{\^e}}1 {à}{{\`a}}1 {ù}{{\`u}}1
           {ô}{{\^o}}1 {î}{{\^i}}1 {ç}{{\c{c}}}1 {É}{{\'E}}1
           {→}{{$\rightarrow$}}1
}

\pagestyle{fancy}
\fancyhf{}
\fancyhead[L]{\small S03 -- Tables de hachage}
\fancyhead[R]{\small Structures de Données Avancées -- M1}
\fancyfoot[C]{\thepage}
\fancyfoot[L]{\small Dr. El Hadji Bassirou TOURÉ}
\fancyfoot[R]{\small 2025--2026}
\renewcommand{\headrulewidth}{0.4pt}
\renewcommand{\footrulewidth}{0.4pt}

\title{%
\vspace{-1cm}
\textbf{\LARGE Structures de Données et}\\
\textbf{\LARGE Algorithmes Avancés}\\[0.8cm]
\textbf{\huge Séance 3 --- Tables de hachage}\\[0.3cm]
\textbf{\Large Comment accéder à une donnée en $\bigO(1)$ ?}\\[0.5cm]
\large Master 1 --- 2025--2026}
\author{Dr. El Hadji Bassirou TOURÉ\\
École Supérieure Polytechnique\\
Département Génie Informatique\\
Université Cheikh Anta Diop de Dakar}
\date{}

\begin{document}

\maketitle
\thispagestyle{empty}

% =====================================================================
% NAVIGATION CONCEPTUELLE
% =====================================================================

\begin{prereqbox}[Navigation conceptuelle]
\textbf{Ce que vous savez déjà :} ADTs Map et Set (S1), complexité $\bigO$/$\Omega$/$\Theta$ (S2). Recherche linéaire $\bigO(n)$, recherche dichotomique $\bigO(\log n)$.\\[4pt]
\textbf{La limite qu'on dépasse aujourd'hui :} \texttt{ArrayList.get(i)} est $\bigO(1)$ \emph{par indice} mais la recherche \emph{par clé} reste $\bigO(n)$. Les listes chaînées ne font pas mieux, les tableaux triés plafonnent à $\bigO(\log n)$ pour \texttt{get} et paient $\bigO(n)$ pour \texttt{put}. Peut-on accéder à une donnée en $\bigO(1)$ \emph{par sa clé} ? Oui --- grâce au hachage.\\[4pt]
\textbf{Ce que vous saurez faire à la fin :}
\begin{enumerate}[nosep]
    \item Expliquer pourquoi aucune structure linéaire ne permet $\bigO(1)$ par clé
    \item Calculer $h(k) = |k.\text{hashCode}()| \bmod m$ pour une clé Java donnée
    \item Dérouler une trace de chaînage séparé et de sondage linéaire
    \item Calculer le facteur de charge $\alpha$ et prévoir les rehashs
    \item Analyser la complexité amortie $\bigO(1)$ moyenne vs $\bigO(n)$ pire cas
\end{enumerate}
\textbf{Ouvre la voie vers :} S4 (arbres de recherche quand l'ordre sur les clés compte).\\[4pt]
\textbf{Mapping CSE 373 :} LEC 7 (Intro to Hashing), LEC 8 (Collision Resolution).\\[4pt]
\textbf{Méthode pédagogique :} Chaque concept majeur est suivi d'un \textbf{mini-exercice d'application} avec sa correction immédiate en vert. Ces exercices sont essentiels --- ne les sautez pas.
\end{prereqbox}

\tableofcontents
\newpage

% =====================================================================
% PARTIE I : LE PROBLÈME ET L'IDÉE-CLÉ
% =====================================================================

\part{Le problème et l'idée-clé}

% =====================================================================
\section{Motivation : au-delà du $\bigO(n)$}
% =====================================================================

\subsection{Le problème concret : un carnet de contacts}

Prenons un cas concret : un \textbf{carnet de contacts téléphoniques}. On y stocke des paires \emph{(nom, numéro)} :

\begin{lstlisting}[style=java]
Map<String, String> carnet = new HashMap<>();
carnet.put("Awa Diop", "+221 77 555 01 23");
carnet.put("Moussa Ndiaye", "+221 76 123 45 67");
carnet.put("Fatou Sall", "+221 77 888 12 34");

String numero = carnet.get("Awa Diop");     // "+221 77 555 01 23"
\end{lstlisting}

Les opérations essentielles sont \texttt{put}, \texttt{get}, \texttt{remove}, \texttt{containsKey}. La question : \textbf{combien coûtent-elles avec les structures déjà vues en S1} ?

\begin{figure}[ht]
\centering
\includegraphics[width=0.95\textwidth]{fig_01_map_complexities.pdf}
\caption{Complexités des implémentations naïves d'un Map. Aucune structure linéaire ne permet d'atteindre $\bigO(1)$ pour les trois opérations.}
\end{figure}

\begin{attentionbox}[Le constat qui motive la séance]
Avec un tableau, une liste chaînée, ou même un tableau trié, la recherche d'un contact dans un carnet de taille $n$ coûte $\bigO(\log n)$ au mieux, $\bigO(n)$ au pire. Pour un service téléphonique d'un million de contacts consulté 10 000 fois par seconde, ce n'est pas acceptable.
\end{attentionbox}

\begin{appliboxQ}[Mini-exercice 1 --- Calcul d'impact]
Une application traite 10\,000 requêtes/seconde sur un carnet de $n = 1\,000\,000$ contacts. Combien d'opérations élémentaires par seconde avec :
\begin{enumerate}[nosep]
    \item Une liste chaînée ($\bigO(n)$) ?
    \item Une table de hachage ($\bigO(1)$) ?
\end{enumerate}
\end{appliboxQ}

\begin{appliboxR}[Réponse]
\begin{enumerate}[nosep]
    \item Liste : $10\,000 \times 1\,000\,000 = 10^{10}$ opérations/seconde. Impossible sur une machine standard ($\approx 10^9$ ops/sec).
    \item Hachage : $10\,000 \times 1 = 10^{4}$ opérations/seconde. Trivial.
\end{enumerate}
Le hachage rend l'application \emph{faisable}, la liste la rend \emph{infaisable}. C'est toute la différence.
\end{appliboxR}

% =====================================================================
\section{Les ADTs associatifs : Set et Map}
% =====================================================================

\begin{infobox}[ADT Map]
Collection de paires $(\text{clé}, \text{valeur})$, chaque clé unique. Opérations : \texttt{put}, \texttt{get}, \texttt{remove}, \texttt{containsKey}, \texttt{size}.
\end{infobox}

\begin{infobox}[ADT Set]
Collection d'éléments distincts. Opérations : \texttt{add}, \texttt{contains}, \texttt{remove}, \texttt{size}.
\end{infobox}

\begin{retenirbox}[Map et Set sont proches]
Un Set peut être vu comme un Map dont on ignore les valeurs. En Java, \texttt{HashSet} est implémenté en interne comme un \texttt{HashMap} sans valeurs. Tout ce qu'on va dire sur le hachage pour Map s'applique à Set.
\end{retenirbox}

% =====================================================================
\section{L'idée-clé : la fonction de hachage}
% =====================================================================

\subsection{Le point de départ : l'accès direct par indice}

En Java, \texttt{arr[i]} coûte $\bigO(1)$ : l'adresse mémoire se calcule par une simple addition.

\textbf{Idée.} Et si on pouvait transformer la clé en un \textbf{indice} valide ? On accéderait directement à la case contenant la valeur, \emph{sans recherche}.

\subsection{Définition}

\begin{infobox}[Fonction de hachage]
Fonction $h : K \to \{0, 1, \ldots, m-1\}$ qui associe à toute clé un indice dans $[0, m-1]$.

\medskip
Deux propriétés fondamentales :
\begin{enumerate}[nosep]
    \item \textbf{Déterminisme} : la même clé produit toujours le même indice.
    \item \textbf{Distribution uniforme} : les clés réparties aussi uniformément que possible.
\end{enumerate}
\end{infobox}

\begin{figure}[ht]
\centering
\includegraphics[width=0.92\textwidth]{fig_02_hash_concept.pdf}
\caption{Chaque clé (nom) est transformée en indice. L'accès à la valeur est direct.}
\end{figure}

\subsection{En Java : \texttt{hashCode()} et compression}

Tout objet Java a une méthode \texttt{hashCode()} qui retourne un \texttt{int}. Pour obtenir un indice dans $[0, m-1]$ :
\[
\boxed{h(\text{clé}) = \bigl|\text{clé.hashCode()}\bigr| \bmod m}
\]

La valeur absolue évite les indices négatifs, le modulo ramène dans $[0, m-1]$.

\begin{theorembox}[Contrat $\texttt{equals} \leftrightarrow \texttt{hashCode}$]
Pour tous objets $a$ et $b$ :
\[
a.\texttt{equals}(b) \;\implies\; a.\texttt{hashCode}() = b.\texttt{hashCode}()
\]
La réciproque n'est \textbf{pas} requise : deux objets différents peuvent avoir le même hashCode (c'est une collision).
\end{theorembox}

\begin{appliboxQ}[Mini-exercice 2 --- Calcul de position]
Soit $m = 11$. Pour chaque clé, calculez $h(\text{clé}) = |\texttt{hashCode()}| \bmod m$.

\medskip
\begin{tabular}{@{}ll@{}}
\texttt{"Awa Diop"} & \texttt{hashCode() = 1234567890} \\
\texttt{"Moussa Ndiaye"} & \texttt{hashCode() = -879654321} \\
\texttt{"Fatou Sall"} & \texttt{hashCode() = 42} \\
\end{tabular}
\end{appliboxQ}

\begin{appliboxR}[Réponse]
\begin{itemize}[nosep]
    \item $h(\texttt{"Awa Diop"}) = |1234567890| \bmod 11 = 1234567890 \bmod 11$. Calcul : $1234567890 = 11 \times 112233444 + 6$, donc $h = 6$.
    \item $h(\texttt{"Moussa Ndiaye"}) = |-879654321| \bmod 11 = 879654321 \bmod 11$. Calcul : $879654321 = 11 \times 79968574 + 7$, donc $h = 7$.
    \item $h(\texttt{"Fatou Sall"}) = 42 \bmod 11 = 9$ (car $42 = 3 \times 11 + 9$).
\end{itemize}
Les trois clés vont donc respectivement en cases [6], [7], [9] --- pas de collision ici.
\end{appliboxR}

% =====================================================================
\section{Le cas idéal : DirectMap sans collisions}
% =====================================================================

Pour comprendre, commençons simplifié : les clés sont déjà des \textbf{entiers bornés} (ex : numéros d'arrêts DDD de 0 à 99). Fonction de hachage = identité : $h(k) = k$.

\begin{figure}[ht]
\centering
\includegraphics[width=0.92\textwidth]{fig_04_directmap.pdf}
\caption{DirectMap : la clé est directement l'indice, pas de collision possible.}
\end{figure}

\subsection{Implémentation}

\begin{lstlisting}[style=java]
public class DirectMap<V> {
    private V[] data;
    private int taille;

    @SuppressWarnings("unchecked")
    public DirectMap(int capacite) {
        data = (V[]) new Object[capacite];
        taille = 0;
    }

    public void put(int cle, V valeur) {
        if (data[cle] == null) taille++;
        data[cle] = valeur;      // (@\textbf{$\bigO(1)$ garanti}@)
    }

    public V get(int cle) {
        return data[cle];         // (@\textbf{$\bigO(1)$ garanti}@)
    }
}
\end{lstlisting}

\begin{retenirbox}[DirectMap : le saint graal, mais…]
Toutes les opérations sont en $\bigO(1)$ \textbf{pire cas}. Problème : ne marche que si le domaine des clés est petit et borné. Pour des noms, c'est inapplicable.
\end{retenirbox}

\begin{attentionbox}[Le dilemme fondamental]
\begin{itemize}[itemsep=2pt]
    \item Si $m$ est \emph{grand} (comme le domaine entier des clés) : gaspillage mémoire.
    \item Si $m$ est \emph{petit} : on compresse beaucoup de clés dans peu de cases $\Rightarrow$ collisions.
\end{itemize}
\end{attentionbox}

% =====================================================================
\section{Les collisions sont inévitables}
% =====================================================================

\begin{theorembox}[Principe des tiroirs]
Si l'on range $n$ objets dans $m$ tiroirs avec $n > m$, au moins un tiroir contient deux objets ou plus.
\end{theorembox}

Appliqué au hachage : dès que $n > m$, deux clés distinctes doivent nécessairement partager un indice.

\begin{figure}[ht]
\centering
\includegraphics[width=0.92\textwidth]{fig_03_pigeonhole.pdf}
\caption{6 clés, 4 cases : au moins deux clés partagent une case --- c'est une collision.}
\end{figure}

\begin{infobox}[Collision]
Deux clés distinctes $k_1 \neq k_2$ telles que $h(k_1) = h(k_2)$.
\end{infobox}

\begin{attentionbox}[Paradoxe des anniversaires]
Avec $m = 365$ cases, il suffit de \textbf{23 clés} pour avoir 50\% de chances de collision. Pour $m = 10\,000$, c'est 119 clés.

\medskip
\textbf{Conclusion.} Il n'est pas question d'\emph{éviter} les collisions. Il faut \textbf{les gérer}.
\end{attentionbox}

\begin{appliboxQ}[Mini-exercice 3 --- Première collision]
Pour $m = 7$, quelles clés parmi ces 5 entrent en collision (ont le même $h$) ?

\medskip
\begin{tabular}{@{}ll@{}}
$k_1$ & \texttt{hashCode = 100} \\
$k_2$ & \texttt{hashCode = 107} \\
$k_3$ & \texttt{hashCode = 50} \\
$k_4$ & \texttt{hashCode = 14} \\
$k_5$ & \texttt{hashCode = 36} \\
\end{tabular}
\end{appliboxQ}

\begin{appliboxR}[Réponse]
Calcul des modulos : $h(k_1) = 100 \bmod 7 = 2$ ; $h(k_2) = 107 \bmod 7 = 2$ ; $h(k_3) = 50 \bmod 7 = 1$ ; $h(k_4) = 14 \bmod 7 = 0$ ; $h(k_5) = 36 \bmod 7 = 1$.

\medskip
\textbf{Collisions :} $(k_1, k_2)$ en case [2], et $(k_3, k_5)$ en case [1]. Les clés $k_1$ et $k_2$ ont même $h$ bien que leurs \texttt{hashCode} diffèrent de 7 (un multiple de $m$).
\end{appliboxR}

\newpage

% =====================================================================
% PARTIE II : RÉSOUDRE LES COLLISIONS
% =====================================================================

\part{Résoudre les collisions}

% =====================================================================
\section{Chaînage séparé (Separate Chaining)}
% =====================================================================

\subsection{Principe}

L'idée la plus simple : si plusieurs clés atterrissent à la même case, on les stocke \textbf{toutes} dans cette case via une \textbf{liste chaînée}. Chaque case = un \emph{seau} (\emph{bucket}) contenant la liste des paires.

\begin{figure}[ht]
\centering
\includegraphics[width=0.92\textwidth]{s3b_fig_01_chaining.pdf}
\caption{Chaînage séparé. Awa Diop, Fatou Sall et Aminata Fall ont toutes $h = 2$ : chaînées en case [2].}
\end{figure}

\subsection{Implémentation}

\begin{lstlisting}[style=java]
public class ChainedHashMap<K, V> {
    private LinkedList<Paire<K, V>>[] table;
    private int m, n;

    private int hash(K cle) {
        return Math.abs(cle.hashCode()) % m;
    }

    public V get(K cle) {
        int i = hash(cle);
        for (Paire<K, V> p : table[i]) {
            if (p.cle.equals(cle)) return p.valeur;
        }
        return null;
    }

    public void put(K cle, V valeur) {
        int i = hash(cle);
        for (Paire<K, V> p : table[i]) {
            if (p.cle.equals(cle)) { p.valeur = valeur; return; }
        }
        table[i].add(new Paire<>(cle, valeur));
        n++;
    }
}
\end{lstlisting}

\begin{appliboxQ}[Mini-exercice 4 --- Trace de chaînage]
Soit une \texttt{ChainedHashMap} avec $m = 5$. On insère les clés (dans cet ordre) :

\medskip
\begin{tabular}{@{}ll@{}}
\texttt{"Awa"}     & $h = 2$ \\
\texttt{"Moussa"}  & $h = 0$ \\
\texttt{"Fatou"}   & $h = 2$ (collision) \\
\texttt{"Cheikh"}  & $h = 4$ \\
\texttt{"Aminata"} & $h = 2$ (collision) \\
\end{tabular}

\medskip
Dessinez l'état final. Combien de comparaisons pour retrouver \texttt{"Aminata"} avec \texttt{get} ?
\end{appliboxQ}

\begin{appliboxR}[Réponse]
État final :
\begin{itemize}[nosep]
    \item [0] : Moussa
    \item [1] : \emph{vide}
    \item [2] : Awa $\to$ Fatou $\to$ Aminata
    \item [3] : \emph{vide}
    \item [4] : Cheikh
\end{itemize}

\medskip
Pour \texttt{get("Aminata")} : on calcule $h = 2$, puis on parcourt la liste [2] : compare avec Awa (non), Fatou (non), Aminata (oui). \textbf{3 comparaisons.}
\end{appliboxR}

\subsection{Analyse du coût moyen}

\begin{theorembox}[Complexité du chaînage]
Facteur de charge $\alpha = n/m$ = nombre moyen de paires par case. Sous hypothèse d'uniformité, une recherche coûte en moyenne :
\[
\boxed{\Theta(1 + \alpha)}
\]
\end{theorembox}

\textbf{Justification.} Calculer $h$ coûte $\bigO(1)$. Parcourir la liste de la case coûte en moyenne $\alpha$. Total : $\Theta(1 + \alpha)$.

\begin{retenirbox}[Conséquence pratique]
Si on maintient $\alpha \leq 0{,}75$ (rehash automatique), le coût moyen est $\Theta(1{,}75) = \Theta(1)$.
\end{retenirbox}

% =====================================================================
\section{Sondage linéaire (Linear Probing)}
% =====================================================================

\subsection{Principe}

Au lieu de listes chaînées, on range tout \textbf{directement dans le tableau}. Si $h(\text{clé})$ est occupée, on essaie $h+1$, puis $h+2$, etc. (modulo $m$), jusqu'à une case libre.

\begin{figure}[ht]
\centering
\includegraphics[width=0.92\textwidth]{s3b_fig_02_probing.pdf}
\caption{Sondage linéaire. Fatou Sall ($h=3$) trouve [3] pris par Awa $\to$ va en [4]. Les déplacements forment un \emph{cluster}.}
\end{figure}

\subsection{Implémentation}

\begin{lstlisting}[style=java]
public V get(K cle) {
    int i = hash(cle);
    while (table[i] != null) {               // (@\textbf{probing}@)
        if (table[i].cle.equals(cle)) return table[i].valeur;
        i = (i + 1) % m;                     // (@\textbf{case suivante}@)
    }
    return null;
}

public void put(K cle, V valeur) {
    int i = hash(cle);
    while (table[i] != null) {
        if (table[i].cle.equals(cle)) { table[i].valeur = valeur; return; }
        i = (i + 1) % m;
    }
    table[i] = new Paire<>(cle, valeur);
    n++;
}
\end{lstlisting}

\begin{appliboxQ}[Mini-exercice 5 --- Trace de sondage]
\texttt{LinearProbingMap} avec $m = 7$. Insertions :

\medskip
\begin{tabular}{@{}ll@{}}
\texttt{"A"} & $h = 3$ \\
\texttt{"B"} & $h = 5$ \\
\texttt{"C"} & $h = 3$ (collision avec A) \\
\texttt{"D"} & $h = 4$ (collision avec C) \\
\end{tabular}

\medskip
Donnez l'état du tableau et les indices finaux. Combien de clés y a-t-il dans le cluster ?
\end{appliboxQ}

\begin{appliboxR}[Réponse]
\begin{itemize}[nosep]
    \item A : va directement en [3]. Tableau : \_ \_ \_ A \_ \_ \_
    \item B : va directement en [5]. Tableau : \_ \_ \_ A \_ B \_
    \item C : $h=3$ occupé. Probing $\to$ [4] libre. Tableau : \_ \_ \_ A C B \_
    \item D : $h=4$ occupé. Probing $\to$ [5] occupé. Probing $\to$ [6] libre. Tableau : \_ \_ \_ A C B D
\end{itemize}

\medskip
\textbf{Cluster de 4 cases :} [3, 4, 5, 6]. Toute nouvelle clé avec $h \in \{3, 4, 5, 6\}$ devra parcourir le cluster entier.
\end{appliboxR}

\subsection{Clustering}

\begin{attentionbox}[Clustering primaire]
Les clés déplacées forment des \textbf{amas} de cases contiguës. Un gros cluster attire encore plus de collisions : toute nouvelle clé tombant \emph{dedans} doit parcourir tout le reste. Les amas tendent à grandir et à fusionner.
\end{attentionbox}

% =====================================================================
\section{Facteur de charge et redimensionnement}
% =====================================================================

\subsection{Le facteur de charge}

\begin{infobox}[Facteur de charge]
\[
\alpha = \frac{n}{m}
\]
Pour le chaînage, $\alpha$ peut dépasser 1. Pour le sondage, $\alpha < 1$ strictement (sinon plus aucune case libre).
\end{infobox}

\subsection{La règle du rehash}

\begin{theorembox}[Règle de redimensionnement]
Quand $\alpha > 0{,}75$ (seuil Java typique) :
\begin{enumerate}[nosep]
    \item Créer un nouveau tableau de taille $m' = 2m$.
    \item \textbf{Recalculer} $h(\text{clé})$ pour chaque paire (car $h$ dépend de $m$).
    \item Réinsérer toutes les paires dans le nouveau tableau.
\end{enumerate}
\end{theorembox}

\begin{figure}[ht]
\centering
\includegraphics[width=0.92\textwidth]{s3b_fig_03_rehash.pdf}
\caption{Redimensionnement : $m = 4 \to 8$. Les positions changent car $h$ dépend de $m$.}
\end{figure}

\begin{appliboxQ}[Mini-exercice 6 --- Calculs de facteur de charge]
\texttt{ChainedHashMap} initialisée à $m = 8$, seuil de rehash $\alpha = 0{,}75$.
\begin{enumerate}[nosep]
    \item Combien de clés distinctes avant le 1\textsuperscript{er} rehash ?
    \item Nouvelle taille $m'$ après ce rehash ? Nouveau $\alpha$ juste après ?
    \item Combien d'insertions supplémentaires avant le 2\textsuperscript{e} rehash ?
\end{enumerate}
\end{appliboxQ}

\begin{appliboxR}[Réponse]
\begin{enumerate}[nosep]
    \item On doit dépasser $\alpha = 0{,}75$ soit $n > 0{,}75 \times 8 = 6$. Donc rehash au moment où $n$ passe de 6 à 7 (après la 7\textsuperscript{e} insertion).
    \item Après rehash : $m' = 16$, $n = 7$, donc $\alpha = 7/16 \approx 0{,}44$.
    \item 2\textsuperscript{e} rehash quand $n > 0{,}75 \times 16 = 12$, soit à la 13\textsuperscript{e} insertion. Insertions supplémentaires depuis le 1\textsuperscript{er} rehash : $13 - 7 = 6$.
\end{enumerate}
\end{appliboxR}

\subsection{Coût amorti}

\begin{attentionbox}[Un rehash coûte cher, mais il est rare]
Un rehash recalcule $n$ hachages : coût $\Theta(n)$. \textbf{Mais} : entre deux rehashs, on a fait au moins $m/2$ insertions. Le coût du rehash est amorti sur ces insertions : chaque insertion paie $\Theta(1)$ pour sa quote-part.

\medskip
\textbf{Conclusion.} Coût amorti d'une insertion : $\bigO(1)$, malgré les rehashs.
\end{attentionbox}

% =====================================================================
\section{Analyse comparative}
% =====================================================================

\subsection{Formules sous uniformité}

\begin{theorembox}[Coût moyen des recherches (sous uniformité)]
\begin{center}
\begin{tabular}{lcc}
\toprule
\textbf{Technique} & \textbf{Succès} & \textbf{Échec} \\
\midrule
Chaînage séparé & $1 + \alpha/2$ & $1 + \alpha$ \\
Sondage linéaire & $\frac{1}{2}(1 + \frac{1}{1-\alpha})$ & $\frac{1}{2}(1 + \frac{1}{(1-\alpha)^2})$ \\
\bottomrule
\end{tabular}
\end{center}
\end{theorembox}

\begin{figure}[ht]
\centering
\includegraphics[width=0.85\textwidth]{s3b_fig_04_cost_alpha.pdf}
\caption{Le chaînage reste linéaire en $\alpha$, le sondage explose dès $\alpha \to 1$.}
\end{figure}

\begin{appliboxQ}[Mini-exercice 7 --- Comparaison numérique]
Quel est le coût moyen d'un \texttt{get} réussi pour :
\begin{enumerate}[nosep]
    \item Chaînage à $\alpha = 0{,}5$ ?
    \item Chaînage à $\alpha = 2$ (pas de rehash) ?
    \item Sondage à $\alpha = 0{,}5$ ?
    \item Sondage à $\alpha = 0{,}9$ ?
\end{enumerate}
\end{appliboxQ}

\begin{appliboxR}[Réponse]
\begin{enumerate}[nosep]
    \item Chaînage, $\alpha = 0{,}5$ : $1 + 0{,}5/2 = 1{,}25$ comparaisons.
    \item Chaînage, $\alpha = 2$ : $1 + 2/2 = 2$ comparaisons. Remarquable : même à $\alpha = 2$, le coût reste raisonnable.
    \item Sondage, $\alpha = 0{,}5$ : $\frac{1}{2}(1 + 1/0{,}5) = \frac{1}{2}(1 + 2) = 1{,}5$ comparaisons.
    \item Sondage, $\alpha = 0{,}9$ : $\frac{1}{2}(1 + 1/0{,}1) = \frac{1}{2}(1 + 10) = 5{,}5$ comparaisons. \textbf{Explosion}.
\end{enumerate}
\end{appliboxR}

\subsection{Tableau récapitulatif}

\begin{center}
\renewcommand{\arraystretch}{1.4}
\begin{tabular}{@{}lccc@{}}
\toprule
 & \textbf{Chaînage} & \textbf{Sondage} & \textbf{Liste/Tableau} \\
\midrule
get / put / remove (moyen) & $\bigO(1)$ & $\bigO(1)$ & $\bigO(n)$ \\
get / put (pire cas)        & $\bigO(n)$ & $\bigO(n)$ & $\bigO(n)$ \\
Mémoire                      & $m$ + chaînons & $m$ & $n$ \\
Cache-friendly               & Non & Oui & --- \\
Robustesse $\alpha$ élevé    & Bonne & Mauvaise & --- \\
\bottomrule
\end{tabular}
\end{center}

\begin{retenirbox}[En pratique]
\textbf{Chaînage} : plus simple, plus robuste. Java utilise le chaînage dans \texttt{HashMap} (avec optimisation rouge-noir si une liste devient trop longue depuis Java 8).

\medskip
\textbf{Sondage} : meilleure utilisation du cache CPU, idéal pour petits objets et $\alpha < 0{,}7$.
\end{retenirbox}

% =====================================================================
% SECTION FINALE : POINTS CLÉS ET TRANSITION
% =====================================================================

\section{Points clés à retenir}

\begin{retenirbox}[Synthèse S03]
\begin{enumerate}
    \item \textbf{Fonction de hachage} $h(k) = |k.\text{hashCode}()| \bmod m$. Déterministe, uniforme.
    \item \textbf{DirectMap} ($h(k) = k$) : $\bigO(1)$ garanti, mais domaine de clés borné.
    \item \textbf{Collisions inévitables} quand $n > m$ (principe des tiroirs).
    \item \textbf{Chaînage séparé} : chaque case = liste chaînée. Coût moyen $\Theta(1 + \alpha)$.
    \item \textbf{Sondage linéaire} : case suivante si occupé. Coût moyen $\Theta(1/(1-\alpha))$, explosif.
    \item \textbf{Facteur de charge} $\alpha = n/m$. Rehash à $\alpha > 0{,}75$ : double $m$, recalcule les positions. Coût amorti $\bigO(1)$.
    \item \textbf{Pire cas} : $\bigO(n)$ partout (adversaire, hachage dégénéré).
    \item \textbf{Contrat Java} : $a.\text{equals}(b) \Rightarrow a.\text{hashCode}() = b.\text{hashCode}()$.
    \item \textbf{Limitation} : pas d'ordre sur les clés $\to$ S4 (arbres).
\end{enumerate}
\end{retenirbox}

\section*{Ce qui vient ensuite : arbres et tas}

\begin{prereqbox}[Vers la séance 4]
\textbf{Ce qu'on sait faire maintenant :} Implémenter une table de hachage complète, gérer le redimensionnement, analyser la complexité amortie. Le hachage donne $\bigO(1)$ moyen.\\[4pt]
\textbf{La limite qu'on va dépasser :} Les tables de hachage n'ont \textbf{aucun ordre}. Impossible de demander « plus petite clé ? », « clés dans $[k_1, k_2]$ », « clés triées ». Ces opérations sont pourtant essentielles (tri alphabétique, plages de dates, intervalles, successeurs\ldots).\\[4pt]
\textbf{En S4 :} Arbres binaires de recherche (BST), arbres AVL avec rotations, tas binaires (priority queues). $\bigO(\log n)$ pour les opérations, \emph{et} ordre maintenu.
\end{prereqbox}

\vfill
\begin{center}
\rule{0.6\textwidth}{0.4pt}\\[6pt]
{\small Dr. El Hadji Bassirou TOURÉ --- ESP/UCAD --- M1 --- 2025--2026}
\end{center}

\end{document}
